\section{Conclusion}
\urgentbr\urgentnote{\textbf{PROPOSED TEXT (Conclusion; paste as \textbackslash{}section\{Conclusion\} here).} Hardware Spectre defenses are prototyped in simulators that can measure their performance but not their security guarantee. \toolname{} closes that gap: a defense is specified as a speculation contract, a protection set and a leakage function, a synthesizer turns the specification into a compact suite of litmus tests that each exhibit a forbidden execution, and an executor runs the suite on the gem5 prototype and checks whether the execution is prevented. On six prototypes, \toolname{} found all five previously known bugs that violate their defense's own guarantee, six new bugs (see Results), in [N]$\times$ fewer tests than random testing on the same defenses. Three limits bound these claims. Each specification is translated from its paper by hand, so a mistake in the translation is a mistake in the test suite, and the protection sets are the least certain part. Enumeration is bounded: Stage~1a is exhaustive up to its instruction bound only for the STT and ReCon models, and SPT's stops at its cap. A hit is a leak only under the window rule and harness we fix (a branch counts as unresolved until its squash lands, and probes never change defense logic), and we call a hit a defect in the prototype only after triage. Because the specification is separate from the tests, checking a new defense or a fix changes only the three definition clauses. Next steps: describe further defenses and speculation sources (for example value prediction), run the executor on RTL, and turn the manual triage that separated checker artifacts from real bugs into automatic checks. \textit{Notes:} the ``six new bugs'' counts the SpecLFB walk order as new (Anna, 2026-11-09); ``[N]$\times$'' is the same placeholder as the abstract and intro (AMuLeT parity run T5 is blocked); the next-steps sentence is a suggestion from the Discussion to-do list, not a result.}